% !TEX root = ../../main.tex

\subsection{Edge Test: Does a Child Depart from Its Parent?}

The child-parent edge test is the first inferential step. For each edge
$(u,c) \in E$, it asks whether child $c$ carries signal beyond its parent $u$.
This is the method's first guard against keeping a split just because the
linkage algorithm happened to create one. In plain terms, the question is:
does this child really move away from the background defined by its parent, or
is the observed difference small enough to be explained by sampling noise?

The raw difference $\theta_c-\theta_u$ is not enough on its own. Some features
are naturally more variable than others, and very small child subtrees are
noisier than large ones. The test therefore begins by converting the raw
child-parent difference into standardized feature scores. The coordinate-wise null is
$H_{0,j}^{\mathrm{edge}}(u,c): \theta_{c,j} = \theta_{u,j}$. Because $c$ is
nested inside $u$, the child and parent estimators are not independent. In the
binary or indicator representation,
\begin{equation}
\mathrm{Var}(\hat{\theta}_{c,j} - \hat{\theta}_{u,j})
=
\theta_{u,j}(1-\theta_{u,j})
\left(\frac{1}{n_c} - \frac{1}{n_u}\right),
\end{equation}
which is the variance expression used in the code. The standardized edge vector
is therefore
\begin{equation}
z_{u \to c,j}^{\mathrm{edge}}
=
\frac{\theta_{c,j} - \theta_{u,j}}
{\sqrt{\theta_{u,j}(1-\theta_{u,j})\left(\frac{1}{n_c} - \frac{1}{n_u}\right)}}.
\end{equation}
This standardization is applied \emph{coordinatewise}. Each feature first gets
its own z-like score before the method combines evidence across features. For a
parent-child edge with $p$ binary or indicator coordinates, the implementation first computes $p$ scalar standardized scores
$z_{u \to c,1}^{\mathrm{edge}},\dots,z_{u \to c,p}^{\mathrm{edge}}$, one for
each feature, then stacks them into the vector
\begin{equation}
z_{u \to c}^{\mathrm{edge}}
=
\left(
z_{u \to c,1}^{\mathrm{edge}},
\dots,
z_{u \to c,p}^{\mathrm{edge}}
\right)^\top
\in \mathbb{R}^p.
\end{equation}
The subtraction term
$\frac{1}{n_c} - \frac{1}{n_u}$ is the signature of the nested tree setting.
In an ordinary two-sample test for independent groups, the variance factor
would look like $\frac{1}{n_1} + \frac{1}{n_2}$. Here the parent already
contains the child, so part of the uncertainty cancels out, which yields the
nested factor $\frac{1}{n_c} - \frac{1}{n_u} > 0$. This is why very small
children need a larger raw deviation $\theta_{c,j} - \theta_{u,j}$ to generate
the same standardized score.

Coordinatewise standardization fixes heteroskedasticity, but it does not solve
the second problem: the edge evidence is high-dimensional and correlated.
Summing $z_{u \to c,j}^2$ directly would overcount redundant coordinates and let
noisy directions dominate just because they contribute many small correlated
terms. The method therefore looks for a smaller set of parent-local directions
that summarize the main variation seen below the parent, and it measures the
edge deviation in that reduced space. For each parent node $u$, the
implementation forms a local matrix $M_u$ whose rows are:
\begin{enumerate}[leftmargin=1.5em]
  \item all descendant leaf vectors $X_i$ for $i \in L(u)$, and
  \item all descendant internal-node distributions $\theta_v$ for proper
  descendants $v$ of $u$.
\end{enumerate}
After removing constant features, let $d_u$ be the number of active
coordinates and let $m_u$ be the number of rows in $M_u$. The code computes the correlation matrix
$C_u \in \mathbb{R}^{d_u \times d_u}$ and its eigen-decomposition
\begin{equation}
C_u = V_u \Lambda_u V_u^\top,
\qquad
\Lambda_u = \mathrm{diag}(\lambda_{u,1},\dots,\lambda_{u,d_u}),
\qquad
\lambda_{u,1} \ge \cdots \ge \lambda_{u,d_u} \ge 0.
\end{equation}
Equivalently, if $\widetilde{M}_u \in \mathbb{R}^{m_u \times d_u}$ denotes the
column-standardized active-coordinate submatrix of $M_u$ after centering each
active feature to mean zero and scaling it to unit variance, then
\begin{equation}
C_u = \frac{1}{m_u}\widetilde{M}_u^\top \widetilde{M}_u,
\end{equation}
up to the standard Pearson-correlation normalization already enforced by the
column standardization. In practice, this matrix identifies the main patterns
of co-variation visible below the parent. When $m_u < d_u$, the implementation
switches to the dual standardized Gram matrix
\begin{equation}
G_u = \frac{1}{d_u}\widetilde{M}_u \widetilde{M}_u^\top,
\end{equation}
whose nonzero eigenvalues encode the same spectral information as those of
$C_u$ up to the normalization convention, so the same local summary is
recovered in a computationally cheaper form.

\paragraph{Why the dual Gram formulation is used.}
The switch from $C_u$ to $G_u$ is not a change of model but a change of
computational viewpoint. When the number of active features $d_u$ is much
larger than the number of local rows $m_u$, directly diagonalizing the
$d_u \times d_u$ feature-space correlation matrix is wasteful: its rank is at
most $m_u$, so most feature-space directions cannot carry additional spectral
information. The dual matrix $G_u$ stores the same row-inner-product structure
in a smaller $m_u \times m_u$ object. This is the same linear-algebra idea
often described as a kernel-trick-style dualization: rather than working with
feature-feature covariance directly, one works with row-row inner products and
then lifts the resulting eigenvectors back into feature space. Concretely, if
$u_{u,i}$ is a unit eigenvector of $G_u$ with eigenvalue $\lambda_{u,i} > 0$,
then the corresponding feature-space direction is recovered by
\begin{equation}
v_{u,i} = \frac{\widetilde{M}_u^\top u_{u,i}}{\sqrt{\lambda_{u,i}}\sqrt{d_u}},
\end{equation}
followed by normalization. The implementation therefore obtains the same
parent-local principal directions while replacing a potentially very large
$d_u \times d_u$ eigendecomposition by a much smaller $m_u \times m_u$ one.

The spectral dimension is the Marchenko--Pastur signal count
\citep{marchenko1967distribution}. With
\begin{equation}
\hat{\sigma}_u^2 = \mathrm{median}\{\lambda_{u,i} : \lambda_{u,i} > 0\},
\qquad
b_u^{\mathrm{MP}} = \hat{\sigma}_u^2 \left(1 + \sqrt{\frac{d_u}{m_u}}\right)^2,
\end{equation}
the projection dimension is
\begin{equation}
k_u^{\mathrm{edge}}
=
\max\left\{
2,\;
\#\{i : \lambda_{u,i} > b_u^{\mathrm{MP}}\}
\right\},
\end{equation}
followed by the implementation cap $k_u^{\mathrm{edge}} \le d_u$. The floor is
$2$.

Let $v_{u,1},\dots,v_{u,k_u^{\mathrm{edge}}}$ denote the leading eigenvectors.
After the coordinatewise standardization step, the vector
$z_{u \to c}^{\mathrm{edge}} \in \mathbb{R}^p$ is projected onto these
parent-local directions. The edge test then rescales those projected
coordinates so that noisy directions do not automatically count more than
stable ones. The resulting test statistic is
\begin{equation}
T_{u \to c}^{\mathrm{edge}}
=
\sum_{i=1}^{k_u^{\mathrm{edge}}}
\frac{(v_{u,i}^\top z_{u \to c}^{\mathrm{edge}})^2}{\lambda_{u,i}}.
\end{equation}
Writing
$w_{u \to c,i} = v_{u,i}^\top z_{u \to c}^{\mathrm{edge}}$ for the projected
coordinate along the $i$th parent-local spectral direction, each summand is
\begin{equation}
\frac{w_{u \to c,i}^2}{\lambda_{u,i}}.
\end{equation}
This is the practical role of the projection-and-whitening step. The projection
rewrites the edge deviation in directions that are actually present in the
subtree below $u$, and the whitening step stops high-variance background
directions from dominating the test merely because they are noisy.
Under the local correlation-model approximation used by the implementation, the
resulting statistic is calibrated against a $\chi^2_{k_u^{\mathrm{edge}}}$
reference law:
\begin{equation}
T_{u \to c}^{\mathrm{edge}} \sim \chi^2_{k_u^{\mathrm{edge}}},
\qquad
p_{u \to c}^{\mathrm{edge}}
=
\Pr\!\left(
\chi^2_{k_u^{\mathrm{edge}}}
\ge
T_{u \to c}^{\mathrm{edge}}
\right).
\end{equation}
This calibration is exact in the pure PCA-whitened case. In the implemented
mixed-basis path, if fewer usable PCA directions are available than the
nominal target dimension, the implementation simply uses the available
parent-local rows rather than padding with unrelated random directions.
\paragraph{Worked child-parent example.}
The worked example below is intended to show the logic of the edge test, not
just the algebra. It starts with a parent-child comparison that is visibly
different by eye, then shows how each formula answers one of the two problems
above: nested variance and correlated high-dimensional evidence. Consider four
binary samples
\begin{equation}
A=(1,1,0,1),\quad
B=(1,1,0,1),\quad
C=(0,0,1,0),\quad
D=(0,0,1,0),
\end{equation}
with tree root $R$ and children $U_1=\{A,B\}$ and $U_2=\{C,D\}$. The node
distributions are
\begin{equation}
\theta_{U_1}=(1,1,0,1),\qquad
\theta_{U_2}=(0,0,1,0),\qquad
\theta_R=(0.5,0.5,0.5,0.5).
\end{equation}
Because the implementation includes proper descendant internal distributions in
the spectral matrix, the root-local matrix is
\begin{equation}
M_R =
\begin{bmatrix}
1&1&0&1\\
1&1&0&1\\
0&0&1&0\\
0&0&1&0\\
1&1&0&1\\
0&0&1&0
\end{bmatrix},
\end{equation}
where the first four rows are the leaves $A,B,C,D$ and the last two rows are
the internal-node distributions $\theta_{U_1}$ and $\theta_{U_2}$. Each active
column has mean $0.5$ and standard deviation $0.5$, so the standardized matrix
is
\begin{equation}
\widetilde{M}_R =
\begin{bmatrix}
1&1&-1&1\\
1&1&-1&1\\
-1&-1&1&-1\\
-1&-1&1&-1\\
1&1&-1&1\\
-1&-1&1&-1
\end{bmatrix}.
\end{equation}
The corresponding parent-local correlation matrix is
\begin{equation}
C_R =
\begin{bmatrix}
1&1&-1&1\\
1&1&-1&1\\
-1&-1&1&-1\\
1&1&-1&1
\end{bmatrix},
\end{equation}
whose eigenvalues are
\begin{equation}
(\lambda_{R,1},\lambda_{R,2},\lambda_{R,3},\lambda_{R,4}) = (4,0,0,0).
\end{equation}
One leading eigenvector is
\begin{equation}
v_{R,1} = (0.5,0.5,-0.5,0.5)^\top.
\end{equation}
Now test the edge $R \to U_1$. Since $n_R=4$ and $n_{U_1}=2$, the nested
variance factor is
\begin{equation}
\frac{1}{n_{U_1}} - \frac{1}{n_R} = \frac{1}{2} - \frac{1}{4} = \frac{1}{4}.
\end{equation}
Because each root coordinate equals $0.5$, every feature has variance
\begin{equation}
0.5(1-0.5)\left(\frac{1}{2}-\frac{1}{4}\right)=0.0625.
\end{equation}
Using the coordinatewise standardization formula above, the edge vector is
\begin{equation}
z_{R \to U_1}^{\mathrm{edge}}
=
\frac{(1,1,0,1)-(0.5,0.5,0.5,0.5)}{\sqrt{0.0625}}
=
(2,2,-2,2)^\top.
\end{equation}
Projecting onto the leading spectral direction gives
\begin{equation}
w_{R \to U_1,1}
=
v_{R,1}^\top z_{R \to U_1}^{\mathrm{edge}}
=
0.5\cdot 2 + 0.5\cdot 2 - 0.5\cdot (-2) + 0.5\cdot 2
=
4.
\end{equation}
Therefore the first whitened contribution to the edge-test statistic is
\begin{equation}
\frac{w_{R \to U_1,1}^2}{\lambda_{R,1}}
=
\frac{4^2}{4}
=
4.
\end{equation}
This example shows the full logic of the child-parent edge test in the same
order as the formulas: standardize the raw difference, express it in the main
parent-level directions, and then down-weight directions that are mostly noise.
In larger real subtrees the implementation can use more than one such
direction, but the core interpretation is already visible in this toy example.
These $p$-values are computed for every edge in the tree.
